Seizure prediction from scalp EEG is difficult to evaluate because patient leakage, class imbalance, overlapping windows, and broad preictal label definitions can produce misleading estimates of performance. This study reports a repository-backed, patient-wise CHB-MIT seizure prediction pipeline and uses completed experiments to examine whether model choice, imbalance handling, or preictal horizon length explains poor discrimination. The frozen dataset contains 10 CHB-MIT patients, 273 EDF recordings, 50 annotated seizures, and 734,796 four-second windows. Each EDF was processed with 18 common EEG channels, 0.5--40 Hz bandpass filtering, 50 Hz notch filtering, 4 s windows with 2 s stride, per-window z-score normalization, a 10 min preictal definition, and a 60 s postictal exclusion interval. EEGNet achieved validation PR-AUC 0.0239 and F1 0.0366; Attention achieved validation PR-AUC 0.0258 and F1 0.0000; and the time-limited Hybrid run achieved validation PR-AUC 0.0238 and F1 0.0461. The label audit found no evidence of mechanical label contamination under the implemented checks, but 47.26\% of positive windows occurred more than five minutes before seizure onset. A controlled preictal horizon ablation over 10, 7.5, 5, and 3 min did not produce a shorter horizon that improved both validation PR-AUC and validation F1 over the 10 min baseline. These results do not support claims of clinical readiness or state-of-the-art performance; instead, they provide an audited negative result showing that simple imbalance handling and global horizon shortening were insufficient under this patient-wise protocol.
